\documentclass[11pt,a4paper]{article}

\usepackage[utf8]{inputenc}
\usepackage[T1]{fontenc}
\usepackage[margin=2.5cm]{geometry}
\usepackage{amsmath,amssymb}
\usepackage{booktabs}
\usepackage[hidelinks]{hyperref}

\title{Differential Privacy in Federated Learning:\\
       Guarantees, Gaps, and the Meaning of $\varepsilon$}
\author{Prepared with LaTeXEditor}
\date{July 2026}

\begin{document}
\maketitle

\begin{abstract}
Federated learning trains a shared model without centralising data, and is
frequently described as privacy-preserving on that basis. It is not. Keeping
data on-device changes who holds the raw records but does not by itself
prevent those records being reconstructed from the updates devices send.
Formal protection requires differential privacy, and composing the two is
harder than either alone. This paper sets out the threat model precisely,
derives the guarantee that federated differential privacy provides, and
examines three gaps between the formal statement and deployed practice: the
difficulty of interpreting the privacy parameter $\varepsilon$, the mismatch
between the unit of privacy and the unit users care about, and the interaction
between privacy noise and the statistical heterogeneity that characterises
federated data. We argue that reported $\varepsilon$ values are widely
misread, that many deployments protect a unit that does not correspond to a
person, and that the accuracy cost of privacy falls unevenly across
participants in ways that aggregate metrics conceal.
\end{abstract}

\tableofcontents

\section{Introduction}
\label{sec:intro}

Federated learning was introduced to train models on data that cannot be
collected \cite{mcmahan2017}. Mobile keyboards, medical records and financial
transactions are all subject to constraints --- regulatory, contractual or
practical --- that make centralising them difficult. The proposal is to send
the model to the data instead: each participant computes an update on local
data, and a server aggregates the updates into a shared model.

The privacy claim usually made for this arrangement is that raw data never
leaves the device. That is true and considerably less protective than it
sounds. The updates a device sends are functions of its data, and functions of
data leak information about data. Gradient inversion attacks reconstruct
recognisable training images from a single update \cite{zhu2019}; membership
inference determines whether a particular record was in the training set
\cite{shokri2017}; and large models memorise and can be induced to emit
verbatim training sequences \cite{carlini2019, carlini2021}.

Data minimisation is therefore necessary but not sufficient. A formal
guarantee requires differential privacy \cite{dwork2006}, which bounds how
much any individual's data can influence the output regardless of what an
adversary knows or subsequently learns. This paper is about what happens when
that formalism meets the federated setting.

\section{Differential privacy}
\label{sec:dp}

\subsection{The definition}

A randomised mechanism $\mathcal{M}$ satisfies $(\varepsilon,
\delta)$-differential privacy if for all adjacent datasets $D, D'$ differing
in one record, and all measurable $S \subseteq \mathrm{range}(\mathcal{M})$,
\begin{equation}
  \Pr[\mathcal{M}(D) \in S]
    \le e^{\varepsilon} \Pr[\mathcal{M}(D') \in S] + \delta.
  \label{eq:dp}
\end{equation}

The definition is a statement about indistinguishability. An adversary
observing the output cannot determine, beyond a factor $e^{\varepsilon}$,
whether any particular record was present. Crucially this holds against an
adversary with arbitrary auxiliary information --- the guarantee does not
depend on assumptions about what else is known, which is precisely the
assumption that anonymisation schemes have historically got wrong.

A full development of the definition and its consequences is given by
\cite{dwork2014}. The parameter $\delta$ admits a failure probability. It should be set well
below $1/n$ for a dataset of $n$ records, since a mechanism that publishes a
randomly chosen record in full satisfies $(0, \delta)$-DP with $\delta = 1/n$
while being catastrophic for the individual concerned.

\subsection{Achieving it}

The standard route is calibrated noise. For a query $f$ with $\ell_2$
sensitivity
\begin{equation}
  \Delta_2 f = \max_{D, D'} \| f(D) - f(D') \|_2,
\end{equation}
the Gaussian mechanism releases $f(D) + \mathcal{N}(0, \sigma^2 I)$ with
$\sigma \ge \Delta_2 f \sqrt{2 \ln(1.25/\delta)} / \varepsilon$.

For gradient-based learning this becomes DP-SGD \cite{abadi2016}: clip each
per-example gradient to a fixed norm $C$, which bounds sensitivity by
construction, then add Gaussian noise:
\begin{equation}
  \tilde{g} = \frac{1}{B}\left(
    \sum_{i=1}^{B} g_i \cdot \min\!\left(1, \frac{C}{\|g_i\|_2}\right)
    + \mathcal{N}(0, \sigma^2 C^2 I) \right).
  \label{eq:dpsgd}
\end{equation}

\subsection{Composition}

Training performs many such releases, and privacy loss accumulates. Basic
composition gives $\varepsilon_{\text{total}} = T\varepsilon$ over $T$ steps,
which is far too pessimistic to be usable --- thousands of steps would exhaust
any reasonable budget. Advanced composition improves this to
$O(\sqrt{T} \varepsilon)$, and the moments accountant \cite{abadi2016} and
Rényi differential privacy \cite{mironov2017} give tighter bounds still by
tracking the full privacy loss distribution rather than its worst case.

These accounting improvements are what make private deep learning feasible at
all. The difference between naive and tight accounting for a realistic
training run is frequently an order of magnitude in $\varepsilon$, which is
the difference between a meaningful guarantee and a vacuous one.

\section{The federated threat model}
\label{sec:threat}

\subsection{Who is the adversary?}

The guarantee depends entirely on who is assumed to be untrusted, and the
federated literature is not always explicit about this.

\paragraph{Central DP.}
The server is trusted to see individual updates and adds noise before
publishing the model. This protects against anyone downstream --- other
participants, users of the released model --- but not against the server
itself. Noise is added once to an aggregate, so utility is comparatively good.

\paragraph{Local DP.}
Each device adds noise before sending, so the server never sees a clean
update. This assumes nothing about the server, which is a much stronger
position, but the cost is severe: noise is added $n$ times and aggregation
averages it down only as $\sqrt{n}$. Local DP at meaningful $\varepsilon$
typically destroys utility for high-dimensional models, though it has been
deployed successfully for simple frequency estimation \cite{erlingsson2014}.

\paragraph{Distributed DP.}
The middle path, and the one most current deployments take. Secure aggregation
\cite{bonawitz2017} lets the server learn only the sum of updates, never any
individual one, using cryptographic masking that cancels when summed. Each
device then adds a fraction of the required noise, so the sum carries the
right total. This achieves central-DP utility under a local-DP-like trust
assumption, at the cost of a protocol that must tolerate devices dropping out
mid-round.

\begin{table}[htbp]
  \centering
  \caption{Trust models and their consequences.}
  \label{tab:trust}
  \begin{tabular}{llll}
    \toprule
    \textbf{Model} & \textbf{Server sees} & \textbf{Noise added} &
      \textbf{Utility} \\
    \midrule
    Central \cite{abadi2016}       & individual updates & once, at server  & best \\
    Distributed \cite{bonawitz2017} & only the sum      & shared across devices & near-central \\
    Local \cite{erlingsson2014}    & noised updates only & once per device & poor \\
    \bottomrule
  \end{tabular}
\end{table}

\subsection{What the updates leak}

It is worth being concrete about why noise is required. Deep leakage from
gradients \cite{zhu2019} shows that for small batches, optimising a synthetic
input to match an observed gradient recovers the original training example
almost exactly. The attack degrades with batch size and with the number of
local steps, which is why federated averaging over many local epochs leaks
less than a single gradient --- but degradation is not prevention, and it
provides no guarantee against an adversary willing to work harder.

\section{Federated DP-SGD}
\label{sec:fedsgd}

\subsection{The unit of privacy}

The most consequential design choice is what constitutes an adjacent dataset
in Equation~\ref{eq:dp}.

\emph{Example-level} privacy protects individual training examples: adjacency
means one example differs. \emph{User-level} privacy protects everything one
participant contributed: adjacency means one participant's entire dataset is
removed \cite{mcmahan2018}. The latter is much stronger, and it is what
federated deployments should provide, because a participant with a thousand
examples receives essentially no protection from an example-level guarantee
--- an adversary learning about all thousand learns about them.

User-level privacy requires clipping the whole model update rather than
per-example gradients, and the sensitivity is then the clipping norm of that
update. This is architecturally simpler than DP-SGD but statistically harder,
because the quantity being clipped is larger and more variable.

\subsection{Amplification by sampling}

Privacy is amplified when the mechanism runs on a random subsample: an
individual not selected in a round incurs no loss from it. Sampling a fraction
$q$ of participants gives roughly $\varepsilon_{\text{eff}} \approx q
\varepsilon$ for small $q$.

Federated deployments have a very large population and sample a small
committee each round, so amplification is substantial and much of the
practical feasibility of federated DP rests on it. The subtlety is that the
formal result requires \emph{uniform random} sampling, and real deployments
sample devices that happen to be charging, idle and on wi-fi. That is not
uniform, it correlates with time zone and socioeconomic factors, and the
amplification claimed is therefore not exactly the amplification obtained.
This gap is acknowledged in the literature \cite{kairouz2021} and rarely
quantified.

\section{What does $\varepsilon$ mean?}
\label{sec:epsilon}

Equation~\ref{eq:dp} is precise, but the number reported alongside it is
routinely over-interpreted.

\paragraph{It is a worst case, not a description.}
$\varepsilon$ bounds the influence of the most influential possible record on
the most revealing possible output. Typical leakage is far below the bound. A
deployment at $\varepsilon = 8$ is not eight times worse than one at
$\varepsilon = 1$ in any behavioural sense, and empirical attack success at
both may be indistinguishable from chance.

\paragraph{It is not comparable across deployments.}
Two systems reporting $\varepsilon = 4$ may differ in the unit of privacy
(example versus user), the adjacency definition, whether the accounting covers
hyperparameter search, and whether $\varepsilon$ is per-round or total. Only
the last of these is usually stated. Comparing the numbers without the
definitions is meaningless.

\paragraph{Large $\varepsilon$ is not nothing.}
Values in the range $\varepsilon = 8$ to $20$ are common in deployed systems
and give weak formal guarantees --- $e^{8} \approx 3000$, so the
indistinguishability bound is nearly vacuous. They nonetheless provide real
protection empirically: the clipping and noise defeat known reconstruction and
memorisation attacks well before they reach the levels the formalism would
require. The honest position is that such deployments have a strong empirical
defence and a weak formal one, and reporting only the $\varepsilon$ obscures
both halves of that.

\section{Heterogeneity}
\label{sec:heterogeneity}

Federated data is not independent and identically distributed. Each
participant's data reflects their own behaviour, so distributions differ
systematically across the population. This is the defining statistical feature
of the setting and it interacts badly with privacy.

\subsection{Effects on optimisation}

Under heterogeneity, local optimisation drives each participant's model toward
their own optimum, and averaging those does not give the global optimum ---
the phenomenon known as client drift. FedProx \cite{li2020} adds a proximal
term penalising divergence from the global model; SCAFFOLD
\cite{karimireddy2020} estimates and corrects the drift directly with control
variates. Both reduce the number of rounds required.

\subsection{Effects on privacy cost}

The interaction that receives least attention is this. Clipping bounds
sensitivity by discarding the portion of an update above the threshold. A
participant whose data is typical produces an update near the population mean
and is clipped lightly. A participant whose data is unusual produces a large,
distinctive update and is clipped heavily --- exactly the participant whose
contribution is most informative, and most in need of protection.

The consequence is that privacy noise degrades accuracy unevenly. Minority
subgroups, whose updates are atypical by definition, lose more of their signal
to clipping and are less well represented in the resulting model. Aggregate
accuracy metrics conceal this entirely: a model can improve on average while
becoming worse for the participants who differ most from the mean.

This is a fairness problem created by a privacy mechanism, and it does not
have a clean solution. Adaptive clipping helps by setting the threshold from a
private quantile estimate rather than a fixed constant, but the underlying
tension --- that protecting outliers requires suppressing exactly what makes
them informative --- is intrinsic.

\begin{table}[htbp]
  \centering
  \caption{Where the accuracy cost of privacy falls.}
  \label{tab:fairness}
  \begin{tabular}{lll}
    \toprule
    \textbf{Participant type} & \textbf{Update magnitude} & \textbf{Effect of clipping} \\
    \midrule
    Typical, large local dataset & near population mean & light \\
    Typical, small local dataset & noisy but centred    & moderate \\
    Atypical / minority subgroup & large and distinctive & severe \\
    \bottomrule
  \end{tabular}
\end{table}

\section{Systems constraints}
\label{sec:systems}

The formal analysis assumes a clean protocol; deployments face conditions that
complicate it.

\paragraph{Dropout.}
Devices leave mid-round --- battery, connectivity, the user picking up the
phone. Secure aggregation must tolerate this, and the recovery protocols add
rounds of communication. Where the noise was distributed across participants,
dropout means the surviving sum carries less noise than the analysis assumed,
so the realised $\varepsilon$ is worse than the reported one unless the
protocol compensates.

\paragraph{Communication.}
Model updates are large and uplink bandwidth is scarce. Compression --- 
quantisation, sparsification, low-rank updates --- is essential and interacts
with privacy: quantisation is itself a randomised mechanism and can be made to
contribute to the privacy guarantee rather than merely degrading utility, an
elegant result that few systems exploit.

\paragraph{Hyperparameter search.}
Choosing the clipping norm and noise multiplier requires experimentation, and
each configuration evaluated on real data consumes privacy budget. Almost no
published deployment accounts for this, so the reported $\varepsilon$
understates the true loss. It is the clearest example of a gap between the
formalism as analysed and the formalism as practised.

\section{Auditing the guarantee}
\label{sec:auditing}

A reported $\varepsilon$ is a claim about an implementation, and
implementations have bugs. Several published privacy analyses have later been
found to overstate protection because the code did not match the analysis ---
gradients clipped after averaging rather than before, noise sampled once and
reused, or an accountant invoked with the wrong sampling rate. The formalism
provides no defence against a mistake in realising it.

Empirical auditing addresses this by measuring privacy loss rather than
deriving it. The standard construction inserts \emph{canaries} --- crafted
records designed to be maximally memorable --- and then tests how well an
adversary distinguishes a model trained with a given canary from one trained
without. The distinguishing advantage yields a lower bound on the true
$\varepsilon$, and if that lower bound exceeds the claimed value, the
implementation is wrong.

The technique has real teeth. Audits have caught implementation errors in
widely used libraries, and the lower bounds obtained are informative in their
own right: for many realistic training configurations the empirical
$\varepsilon$ is far below the analytical one, which quantifies the looseness
discussed in Section~\ref{sec:epsilon}. A deployment reporting an analytical
$\varepsilon = 10$ and an audited lower bound of $1.5$ is telling a more
complete story than one reporting either alone.

\paragraph{Limitations.}
Auditing produces a lower bound, so it can falsify a claim but never confirm
one. A canary that fails to be extracted shows the adversary tried and failed,
not that a better adversary would. Audits also require training many models,
which is expensive at realistic scale, so they are typically performed on
reduced configurations whose relationship to the deployed one is assumed
rather than demonstrated.

\section{Alternatives and complements}
\label{sec:alternatives}

Differential privacy is not the only technology in this area, and the others
address different parts of the problem.

\paragraph{Homomorphic encryption.}
Computation on encrypted data, so the server aggregates without decrypting.
This removes the need to trust the server with individual updates and provides
a cryptographic rather than statistical guarantee. It does not, however,
protect the \emph{output}: a model trained on encrypted data and then
decrypted memorises exactly as much as one trained in the clear. Homomorphic
encryption is a substitute for secure aggregation, not for differential
privacy, and the two solve orthogonal problems. Its practical obstacle remains
cost, though schemes tailored to the aggregation pattern have narrowed the gap
considerably.

\paragraph{Trusted execution environments.}
Hardware enclaves offer a pragmatic alternative: the server computes on plain
data inside a region it cannot inspect, attested by the processor. This is
much faster than cryptographic approaches and shifts the trust to the silicon
vendor. A steady stream of side-channel vulnerabilities has shown that trust
to be conditional, and enclaves are best regarded as raising the cost of
attack rather than eliminating it.

\paragraph{Synthetic data.}
Generating a private synthetic dataset and training on it freely is
attractive, since the privacy cost is paid once and downstream use is
unconstrained. The difficulty is that synthetic data generated under a strict
privacy budget preserves marginal distributions well and higher-order
structure poorly, and models trained on it underperform in ways that are hard
to predict from summary statistics. The approach works best where the
downstream task is known in advance, so the generator can be tuned to preserve
what that task needs.

\section{Regulation and the gap it leaves}
\label{sec:regulation}

The legal frameworks that motivate much of this work do not map cleanly onto
the formalism, and the mismatch has practical consequences.

Data protection regimes generally distinguish personal data, subject to
extensive obligations, from anonymous data, which is largely unregulated. The
question of where differentially private output falls is genuinely unsettled.
A model trained with a small $\varepsilon$ is intuitively anonymous, but the
regimes were drafted around a binary notion of identifiability that does not
accommodate a continuous parameter, and no widely accepted threshold exists.

Two consequences follow. First, organisations have limited incentive to choose
a small $\varepsilon$ over a large one, since neither reliably changes their
regulatory position; the choice is made on internal risk judgement instead.
Second, the right to erasure sits awkwardly with trained models. Removing a
participant's data from a training set does not remove its influence from a
model already trained, and retraining is often infeasible. Machine unlearning
addresses this directly, and differential privacy addresses it indirectly:
a guarantee that no individual much influenced the model is close to a
guarantee that removing them would change little.

\paragraph{What would help.}
Guidance that named an $\varepsilon$ threshold, however arbitrary, would
change practice more than any technical advance --- it would convert a
judgement call into a compliance target and give organisations a reason to
prefer stronger settings. The absence of such guidance is currently the main
reason deployments cluster at the weak end of the range.

\section{Practical recommendations}
\label{sec:recommendations}

For practitioners deploying federated learning with privacy, the preceding
analysis suggests the following.

\paragraph{State the adjacency definition.}
Report whether the guarantee is example-level or user-level, and if user-level,
what bounds a single user's contribution. Without this the $\varepsilon$ is
uninterpretable.

\paragraph{Account for everything.}
Include hyperparameter search, model selection, and any pilot runs on real
data in the budget. If this is impractical, say so explicitly rather than
reporting a figure that covers only the final run.

\paragraph{Audit the implementation.}
Run a canary-based audit at reduced scale before deployment. It is the only
practical check that the code implements the analysis.

\paragraph{Report distributional utility.}
Publish accuracy by subgroup, not only in aggregate. The clipping analysis of
Section~\ref{sec:heterogeneity} predicts that the cost falls on atypical
participants, and only disaggregated metrics reveal whether it has.

\paragraph{Prefer distributed DP.}
Secure aggregation with distributed noise gives central-DP utility without
requiring participants to trust the server with their individual updates. The
protocol complexity is real but the trust reduction is substantial, and mature
implementations now exist.

\section{Cross-silo and cross-device: two different problems}
\label{sec:silo}

The literature treats federated learning as one setting; it is two, and the
privacy analysis differs substantially between them.

\emph{Cross-device} federation involves an enormous population of unreliable
participants each holding little data --- the mobile keyboard case. Sampling
amplification is strong because the sampling rate is tiny, participants are
anonymous and interchangeable, and no individual round depends on any
particular device. Dropout is expected and designed for.

\emph{Cross-silo} federation involves a small number of reliable institutional
participants each holding a great deal of data --- hospitals collaborating on a
diagnostic model, or banks on fraud detection. Here almost every assumption
inverts. With ten participants there is no meaningful sampling amplification.
Participants are identifiable and persistent, so an update is attributable
even before its contents are considered. And the unit of privacy becomes
ambiguous in a way it is not on a phone: protecting the \emph{hospital} is
usually not what matters, since the hospital is a party to the collaboration;
protecting each \emph{patient} within it is, and that requires example-level
or group-level accounting inside each silo as well as across them.

The consequence is that cross-silo deployments generally cannot achieve strong
user-level DP by the mechanisms of Section~\ref{sec:fedsgd}, and typically
rely more heavily on secure aggregation, contractual controls and access
restriction, with differential privacy applied at a weaker setting as one
layer among several. Papers reporting federated DP results should state which
regime they address, because a technique that works in one may be inapplicable
in the other.

\section{Discussion}
\label{sec:discussion}

Three conclusions follow from the preceding sections.

First, federated learning without differential privacy should not be described
as privacy-preserving. It provides data minimisation, which is valuable and
distinct. The reconstruction results \cite{zhu2019} are sufficient to
establish that raw updates leak, and any claim of privacy needs to say what
mechanism provides it.

Second, the unit of privacy must be stated. User-level guarantees are what the
setting calls for and example-level guarantees are what are sometimes
reported, and the difference is not a technicality --- for a participant
contributing many correlated records, an example-level guarantee is close to
no guarantee at all.

Third, aggregate utility metrics are inadequate for evaluating private
federated systems. The cost of privacy is borne disproportionately by
atypical participants, so a system should report per-subgroup performance.
Reporting mean accuracy alongside an $\varepsilon$ describes neither the
protection nor its distribution.

\section{Conclusion}
\label{sec:conclusion}

Differential privacy provides the only guarantee in this area that survives an
adversary with unknown auxiliary information, and combining it with federated
learning and secure aggregation gives a deployable system with a real formal
statement attached. That is a genuine achievement and it is deployed at scale.

The gaps are in interpretation rather than mechanism. $\varepsilon$ is a
worst-case bound routinely read as a description of typical leakage; it is
compared across deployments whose definitions differ; the sampling
amplification that makes the numbers work assumes uniformity that device
availability violates; hyperparameter search is uncounted; and the accuracy
cost falls hardest on the participants least represented in the data.

None of these invalidates the approach. They do mean that a reported
$(\varepsilon, \delta)$ pair is the beginning of a privacy analysis rather
than its conclusion, and that the field would be better served by reporting
the adjacency definition, the accounting scope, and the distribution of the
utility cost alongside the parameter.

\begin{thebibliography}{99}

\bibitem{abadi2016}
M.~Abadi, A.~Chu, I.~Goodfellow, H.~B. McMahan, I.~Mironov, K.~Talwar, and
  L.~Zhang.
\newblock Deep learning with differential privacy.
\newblock In \emph{ACM CCS}, 2016.

\bibitem{bonawitz2017}
K.~Bonawitz, V.~Ivanov, B.~Kreuter, A.~Marcedone, H.~B. McMahan, S.~Patel,
  D.~Ramage, A.~Segal, and K.~Seth.
\newblock Practical secure aggregation for privacy-preserving machine
  learning.
\newblock In \emph{ACM CCS}, 2017.

\bibitem{carlini2019}
N.~Carlini, C.~Liu, Ú.~Erlingsson, J.~Kos, and D.~Song.
\newblock The secret sharer: Evaluating and testing unintended memorization in
  neural networks.
\newblock In \emph{USENIX Security}, 2019.

\bibitem{carlini2021}
N.~Carlini et~al.
\newblock Extracting training data from large language models.
\newblock In \emph{USENIX Security}, 2021.

\bibitem{dwork2006}
C.~Dwork, F.~McSherry, K.~Nissim, and A.~Smith.
\newblock Calibrating noise to sensitivity in private data analysis.
\newblock In \emph{TCC}, 2006.

\bibitem{dwork2014}
C.~Dwork and A.~Roth.
\newblock The algorithmic foundations of differential privacy.
\newblock \emph{Foundations and Trends in Theoretical Computer Science},
  9(3--4):211--407, 2014.

\bibitem{erlingsson2014}
Ú.~Erlingsson, V.~Pihur, and A.~Korolova.
\newblock RAPPOR: Randomized aggregatable privacy-preserving ordinal response.
\newblock In \emph{ACM CCS}, 2014.

\bibitem{kairouz2021}
P.~Kairouz et~al.
\newblock Advances and open problems in federated learning.
\newblock \emph{Foundations and Trends in Machine Learning}, 14(1--2):1--210,
  2021.

\bibitem{karimireddy2020}
S.~P. Karimireddy, S.~Kale, M.~Mohri, S.~J. Reddi, S.~U. Stich, and A.~T.
  Suresh.
\newblock SCAFFOLD: Stochastic controlled averaging for federated learning.
\newblock In \emph{ICML}, 2020.

\bibitem{li2020}
T.~Li, A.~K. Sahu, M.~Zaheer, M.~Sanjabi, A.~Talwalkar, and V.~Smith.
\newblock Federated optimization in heterogeneous networks.
\newblock In \emph{MLSys}, 2020.

\bibitem{mcmahan2017}
H.~B. McMahan, E.~Moore, D.~Ramage, S.~Hampson, and B.~A. y~Arcas.
\newblock Communication-efficient learning of deep networks from decentralized
  data.
\newblock In \emph{AISTATS}, 2017.

\bibitem{mcmahan2018}
H.~B. McMahan, D.~Ramage, K.~Talwar, and L.~Zhang.
\newblock Learning differentially private recurrent language models.
\newblock In \emph{ICLR}, 2018.

\bibitem{mironov2017}
I.~Mironov.
\newblock Rényi differential privacy.
\newblock In \emph{IEEE Computer Security Foundations Symposium}, 2017.

\bibitem{shokri2017}
R.~Shokri, M.~Stronati, C.~Song, and V.~Shmatikov.
\newblock Membership inference attacks against machine learning models.
\newblock In \emph{IEEE S\&P}, 2017.

\bibitem{zhu2019}
L.~Zhu, Z.~Liu, and S.~Han.
\newblock Deep leakage from gradients.
\newblock In \emph{NeurIPS}, 2019.

\end{thebibliography}

\end{document}
